\scanpage{110}
% archive sheet number: 109
a) Donnons nous un idéal $R\subset B$. Alors

Décomposition de $\deleted{\unsure{A}}B/R$ en produit de deux anneaux

$\simeq$ Donnée de deux idéaux $J,J'$ de $A$ tels que
\[
JJ'=R,\quad J+J'=A\qquad [\text{Car }J+J'=A\Rightarrow JJ'=J\cap J']
\]

b) Supposons $B=A[T]$, \inserted{\deleted{$A$ \unsure{local donc}} $A$ connexe} $R=(P)$, $P$ polynôme
\uline{unitaire}, Alors tte décomposition de $R$ en
produit de deux idéaux $J,J'$ tels que $J+J'=B$,
est uniquement de la forme $(\deleted{\unsure{R}}P)=(Q)(Q')$, où
$Q,Q'$ sont des polynômes \uline{unitaires} tels que
$P=QQ'$.

Si elle est \unsure{donnée|dominée} par une telle déc,
alors $P$ est unique à une unité près,
et étant supposé \unsure{unifaire|unitaire} est unique
\deleted{[N.B. tout} [petite vérification \ldots.]. \deleted{On
\unsure{prouvera que}} De l'unicité résulte que
l'on peut supposer $\operatorname{Spec} A$ \uline{local}. Alors pour
l'existence, on note que c'est vrai sur les
